% !TeX root = ../../Book3.tex
\chapter{Poincaré 度量与双曲几何}
\label{b3:ch:13}
\BookChapterNumber{b3:ch:13}
\begin{chapterabstract}
Schwarz--Pick 不等式中的密度因子可以用来测量路径，由此得到圆盘上的 Poincaré 距离。本章从长度下确界出发，证明距离公式、测地线的唯一性、完备性及曲率为 \(-1\) 的归一化，再把这些结论转写到上半平面。一般双曲曲面通过圆盘覆盖取得同样的局部几何；不同提升之间的选择则决定底曲面的全局距离。全纯压缩、密度单调性、商曲面的最短路和尖端计算将说明，局部等距、全局保距、有限面积及完备性各自控制什么。
\end{chapterabstract}

\begin{learninggoals}
\begin{itemize}
\item\label{b3:item:13-goal-1}从路径长度推出圆盘距离，找出达到下确界的路径，并解释测地线为何是正交圆弧或直径。
\item\label{b3:item:13-goal-2}核对曲率为 \(-1\) 的密度、距离和面积常数，在圆盘与半平面之间计算，并区分理想边界与度量完备化。
\item\label{b3:item:13-goal-3}借覆盖提升建立曲面上的 Poincaré 度量与 Schwarz--Pick 不等式，准确说明导数等号为何给覆盖。
\item\label{b3:item:13-goal-4}用甲板轨道比较全局距离，理解紧闭球、单射半径和最短路的关系，并将这些性质用于映射族与曲面端部。
\end{itemize}
\end{learninggoals}

圆盘自同构在欧氏坐标中的导数可以大于1，也可以小于1。第六章的 Schwarz--Pick 估计说明，若同时比较源点和像点处的密度 \(2/(1-|z|^2)\)，这种伸缩恰好抵消。把密度沿曲线积分以后，自同构保持每条路径的长度，一般全纯映射则使长度不增。问题随之从两个函数值的分式估计，转为哪些路径最短、哪些点列能够收敛，以及绕孔是否能产生更短的连接。

\ProseParagraphBreak
\secref{b3:sec:1301}先把长度元、路径长度与两点距离分开定义，再从径向积分求距离公式。\secref{b3:sec:1302}沿着等号条件寻找测地线，并检查完备性、圆盘面积和曲率；\secref{b3:sec:1303}换到半平面，使竖直运动、水平运动和趋近理想边界都能直接计算。局部的欧氏直觉仍然有用，但不能把“欧氏距离很近”直接等同于“双曲距离很近”，也不能把一条便于积分的路径当作最短路。

\ProseParagraphBreak
最后的\secref{b3:sec:1304}使用第十一章的路径提升、映射提升和甲板群，把圆盘几何传到一般曲面。密度是局部对象，覆盖导数和自同构不变性足以保证它下降；距离却要在所有连接路径中取下确界，因此必须比较同一纤维的全部提升。这个差别解释了为何覆盖处处局部等距，仍可能缩短给定两点的距离，且底曲面的最短路可能不唯一。

\ProseParagraphBreak
这些问题也为曲率赋予具体意义：圆盘边缘的密度增长使边界距离无穷远，而曲率公式描述这一长度单位的局部变化。\cref{b3:thm:13-cover-model-input}给出曲面具有圆盘覆盖的条件；借助这个模型，圆盘中的计算可以用于更一般的曲面。第十四章将进一步研究负曲率度量的比较及正规族问题。

% !TeX root = ../../Book3.tex
\section{圆盘上的 Poincaré 度量}
\label{b3:sec:1301}

Schwarz--Pick 的分母给出圆盘内的局部长度单位。把这个单位沿曲线积分，再在曲线中取下确界，就得到距离；分式三角不等式则变成测地线上的普通加法。

\subsection{度量元}
\label{b3:subsec:130101}

圆盘的\term[poincare-duliang]{Poincaré 度量}[Poincaré metric]长度元为
\[
 \mathrm{d}s=\lambda_{\mathbb D}(z)|\mathrm{d}z|,
 \qquad \lambda_{\mathbb D}(z)=\frac2{1-|z|^2}.
\]
分段 \(C^1\) 路径的长度为 \(L(\gamma)=\int\lambda_{\mathbb D}(\gamma(t))|\gamma'(t)|\,\mathrm{d}t\)。本卷固定曲率为 \(-1\) 的归一化，不能删掉因子2而仍沿用下面的曲率和距离式。

\ProseParagraphBreak
密度在内部光滑且正，趋近边界时发散。内部小紧盘上它与欧氏密度有统一双边界，所以给相同局部拓扑；发散则把有限欧氏距离的边界推到无限内在距离。

\ProseParagraphBreak
度量元不是点之间的距离。它先给分段光滑路径长度 \(\ell(\gamma)=\int\lambda(\gamma)|\gamma'|\,\mathrm{d}t\)，再对连接两点的全部路径取下确界。因为 \(\lambda\) 在每个紧圆盘上有正上下界，局部长度与 Euclid 长度可比较；因此不同点的距离为正，且度量拓扑与原来拓扑相同。对圆盘中远近的比较不能仅把两个端点的密度平均，因为最短路径可能绕开密度大的区域。

\ProseParagraphBreak
本章归一化曲率为 \(-1\)，故密度是 \(2/(1-|z|^2)\)。也有文献使用 \(1/(1-|z|^2)\)，所有长度减半，曲率成为 \(-4\)。换常数并不改变测地线形状，但会改变距离、面积和闭曲线长度中的系数。后面圆环核心长度与模量的乘积为 \(\pi\)，也是这一归一化的结果。

\subsection{Möbius 不变性}
\label{b3:subsec:130102}

圆盘自同构满足 \(\lambda_{\mathbb D}(Tz)|T'(z)|=\lambda_{\mathbb D}(z)\)，已由第6章的导数恒等式得到。换元使每条路径长度不变，因此也保持路径距离。这个不变性把任意点的计算归到0。

\ProseParagraphBreak
对一般全纯圆盘映射只有不等式 \(\lambda_{\mathbb D}(f(z))|f'(z)|\le\lambda_{\mathbb D}(z)\)。于是每条像路径长度不超过原路径长度。若在一点等号且映射非常值，Schwarz--Pick 刚性使它成为自同构；长度不变的全局对称由局部导数等号识别。

\subsection{Poincaré 距离}
\label{b3:subsec:130103}

\begin{theorem}\label{b3:thm:13-disk-distance}
Poincaré 路径距离为
\[
 \mathrm d_{\mathbb D}(z,w)=2\operatorname{arctanh}\delta(z,w)
 =\log\frac{1+\delta(z,w)}{1-\delta(z,w)}.
\]
\end{theorem}
\begin{proof}
先令一端为0，另一端模为 \(r\)。对任意连接路径，记径向函数 \(u(t)=|\gamma(t)|\)，几乎处处 \(|u'|\le|\gamma'|\)。故其长度至少为 \(\int 2|u'|/(1-u^2)\,\mathrm{d}t\)，又至少为原函数 \(2\operatorname{arctanh}u\) 的端点差 \(2\operatorname{arctanh}r\)。径向线段达到此值。再用自同构把 \(w\) 移到0，像点模为 \(\delta(z,w)\)，不变性给公式。
\end{proof}
密度正使不同点距离正；路径拼接给三角不等式，反向给对称。第6章分式三角式经反双曲正切后也直接验证相同结论。

\begin{proposition}[距离的双曲余弦形式]\label{b3:prop:13-cosh-distance}
对 \(z,w\in\mathbb D\)，有
\[
 \cosh d(z,w)=1+\frac{2|z-w|^2}{(1-|z|^2)(1-|w|^2)}.
\]
若两条从0出发的径向线夹角为 \(\theta\)，端点到0的距离为 \(a,b\)，端点间距离为 \(c\)，则
\[
 \cosh c=\cosh a\cosh b-\sinh a\sinh b\cos\theta.
\]
\end{proposition}
\begin{proof}
令 \(q=\delta(z,w)\)，\(d=2\operatorname{artanh}q\)，故 \(\cosh d=(1+q^2)/(1-q^2)\)。再用
\[
 |1-\bar wz|^2-|z-w|^2=(1-|z|^2)(1-|w|^2)
\]
化简即得第一式。第二式取 \(z=r\)、\(w=s e^{\mathrm{i}\theta}\)，则分子为 \(r^2+s^2-2rs\cos\theta\)。由 \(r=\tanh(a/2)\)、\(s=\tanh(b/2)\)，\(\cosh a=(1+r^2)/(1-r^2)\)、\(\sinh a=2r/(1-r^2)\)，代入后逐项整理得所写公式。
\end{proof}
因共形度量保持夹角，\(\theta\) 也是该三角形在0的内角。把任一顶点用圆盘自同构送到0，便得到一般双曲三角形在该顶点的余弦关系。这个推导只用已经建立的距离公式，无需先接受一般曲面的三角公式。

\subsection{Schwarz–Pick 压缩性}
\label{b3:subsec:130104}

\begin{proposition}\label{b3:prop:13-disk-contraction}
圆盘全纯自映射压缩 Poincaré 距离。不同两点之间距离取等，或一个点的密度导数取等，都迫使映射是自同构。
\end{proposition}
\begin{proof}
沿任意源路径积分 Schwarz--Pick 导数不等式，得到像长度不大于源长度。目标距离不大于像长度，再对源路径取下确界。等号的两点形式也可直接对严格递增的 \(2\operatorname{arctanh}\) 使用伪双曲等号刚性；导数形式在第6章已证明。
\end{proof}
距离压缩比较的是域的内在单位。\(z\mapsto(z+a)/(1+az)\) 可有欧氏导数大于1，但它处处保持这里的长度。

\begin{proposition}[固定像值的紧致控制]\label{b3:prop:13-anchored-maps}
设 \(X\) 为圆盘型曲面，\(K\subset X\) 紧。全纯映射族 \(f:\mathbb D\to X\)、\(f(0)\in K\)，在每个闭子盘的像都包含于 \(X\) 的一个紧集，并且每个序列有局部一致收敛到全纯映射的子列。
\end{proposition}
\begin{proof}
第四节的定理\ref{b3:thm:13-hyperbolic-schwarz-pick}给 \(d_X(f(z),f(0))\le d_{\mathbb D}(z,0)\)。固定 \(|z|\le r<1\)，右边有统一上界。用有限个球覆盖 \(K\)，并扩大半径，像落在有限个紧闭球的并中；紧闭球性质由本章后面独立的覆盖距离证明给出。同样的压缩使族在每个子盘等度连续。在每个子盘选可数稠密点，利用紧目标逐点抽取并作对角子列。给定误差，等度连续性提供统一小邻域，再用有限网格覆盖子盘；有限网格点的收敛和三角不等式使子列在整个子盘一致 Cauchy。对可数圆盘穷竭再作对角抽取，得到局部一致极限 \(F\)。在任一点附近，把 \(F\) 的像限制到目标小坐标盘；一致收敛使子列充分晚也落在该盘中，其坐标表达的全纯极限由 Weierstrass 定理全纯。因此 \(F\) 是全纯曲面映射。
\end{proof}
固定0的像落在紧集的条件阻止整族向曲面端部逃走。在三穿孔球面上，它给省略值函数的一种紧致控制，但若基点像趋向穿孔，就可能只有球面值的退化极限。第十四章采用球面正规定义，能同时容纳这类退化；本命题则给非退化的曲面值极限。这两种紧致性陈述的目标不同，使用前应当检查需要哪一种极限。

% !TeX root = ../../Book3.tex
\section{双曲测地线}
\label{b3:sec:1302}

圆盘的最短路是直径的自同构像。完备性来自边界的无限距离，而常负曲率则描述同一局部尺度的微分几何。

\subsection{测地线}
\label{b3:subsec:130201}

\begin{theorem}\label{b3:thm:13-disk-geodesics}
圆盘中任意不同两点由唯一的测地线段连接，忽略保持方向的重参数化。测地线是与单位圆正交的圆弧或直径。
\end{theorem}
\begin{proof}
以自同构把一端移到0。距离证明中的两个不等式同时取等，要求路径速度几乎处处径向且半径单调，故只有径向线段能最短。撤销自同构后得到唯一线段。Möbius 变换保持广义圆与夹角，直径与单位圆正交，所以像为正交圆弧。反向每条这样的圆弧可由一个圆盘自同构送成直径，故确为测地线。
\end{proof}
沿直径可用弧长参数 \(r=\tanh(s/2)\)，\(-\infty<s<\infty\)，说明完整测地线的两端都在圆周，而不会在有限弧长到达边界。

\subsection{完备性}
\label{b3:subsec:130202}

\begin{theorem}\label{b3:thm:13-disk-complete}
圆盘 Poincaré 度量完备，且闭有界距离球为紧集。
\end{theorem}
\begin{proof}
0点距离公式使半径 \(M\) 的闭球为欧氏闭盘 \(|z|\le\tanh(M/2)<1\)，故紧；其他中心由自同构传回。Cauchy 列到某个固定点的距离有界，因此落在这样的紧球。密度至少2，故它也是欧氏 Cauchy 列，并有盘内极限；在包含该球的略大紧盘上密度有上界，直线段估计使它在 Poincaré 距离也趋于这个极限。
\end{proof}
完备不等于底空间欧氏闭合。开圆盘在欧氏空间不完备，在这个长度单位中却完备。度量改变了 Cauchy 列的条件，使逃向边界的点列不再 Cauchy。

\begin{proposition}[双曲圆的长度与面积]\label{b3:prop:13-ball-area}
双曲半径 \(R\) 的圆盘面积与边界长度分别为
\[
 A(R)=2\pi(\cosh R-1)=4\pi\sinh^2(R/2),\qquad
 L(R)=2\pi\sinh R.
\]
\end{proposition}
\begin{proof}
由自同构不变性，只需中心0。双曲半径 \(R\) 对应 Euclid 半径 \(r=\tanh(R/2)\)。共形面积元是 \(\lambda^2\,\mathrm{d}A\)，故
\[
 A=\int_0^r\frac{8\pi t}{(1-t^2)^2}\,\mathrm{d}t
   =\frac{4\pi r^2}{1-r^2}=4\pi\sinh^2(R/2).
\]
圆周上密度恒为 \(2/(1-r^2)\)，长度为 \(4\pi r/(1-r^2)=2\pi\sinh R\)。
\end{proof}
小半径时 \(A(R)\sim\pi R^2\)、\(L(R)\sim2\pi R\)，符合局部 Euclid 行为；大半径时两者均指数增长。完整双曲盘的总面积无穷，而每个有限半径闭球紧。这两项性质相容，完备性不意味着总面积有限。

\ProseParagraphBreak
闭球紧也可从距离公式直接看。以0为中心的闭球是 \(|z|\le\tanh(R/2)<1\) 的 Euclid 闭盘；任意中心由自同构化到0。每条有界距离的序列因此有收敛子列，这比仅有 Cauchy 完备性更强。后面甲板轨道的距离下确界能够达到，正是依赖闭球紧与轨道离散共同作用。

\subsection{常负曲率}
\label{b3:subsec:130203}

对共形长度元 \(\lambda|\mathrm{d}z|\)，Gaussian 曲率计算式为
\[
 K=-\lambda^{-2}\Delta\log\lambda.
\]
令 \(u=\log\lambda\)，度量张量为 \(e^{2u}(\mathrm{d}x^2+\mathrm{d}y^2)\)。其 Christoffel 系数为 \(\Gamma^1_{11}=u_x,\Gamma^1_{12}=u_y,\Gamma^1_{22}=-u_x\)，另三项由对称交换坐标得到；代入曲率张量后，一次导数的平方项相消，留下 \(-e^{-2u}(u_{xx}+u_{yy})\)。因此上述式子没有使用统一化。

\ProseParagraphBreak
对圆盘 \(u=\log2-\log(1-|z|^2)\)，直接计算 \(\Delta u=4/(1-|z|^2)^2=\lambda^2\)，所以 \(K=-1\)。常数密度的欧氏度量曲率0，不能通过整体缩放变成负曲率；位置变化的密度是决定性因素。

\subsection{等距变换}
\label{b3:subsec:130204}

\begin{proposition}\label{b3:prop:13-disk-isometries}
保定向光滑 Poincaré 等距变换恰为圆盘 Möbius 自同构；反定向光滑等距变换是它们与复共轭的复合。
\end{proposition}
\begin{proof}
等距条件使导数在每点为欧氏相似变换。保定向时导数满足 Cauchy--Riemann 方程，故变换全纯；它及逆均全纯，因此由圆盘自同构分类为 Möbius。反定向时先复合共轭转成保定向。反向由密度不变性和共轭不变性得到。
\end{proof}
这里给出光滑 Riemannian 等距变换的分类，没有把一般度量空间等距映射的自动光滑性作为未证明前提。

\ProseParagraphBreak
径向线虽都经过0，通常的 Euclid 直线段却不一定是双曲测地线。两点若不在同一直径上，连接它们的最短曲线是与单位圆正交的圆弧。圆弧可以通过将第一点送到0的 Möbius 变换取得，再把径向段拉回。Möbius 变换保持圆及直线的广义圆，也保持交角，所以得到正交性。弧的唯一性来自拉回后径向最短路的唯一性，而不是仅凭图中能画出一条圆弧。

\ProseParagraphBreak
圆盘自同构可能没有内部不动点，这不妨碍它为等距变换。抛物型和双曲型的不动点在边界，而边界处距离无穷，不属于度量空间。椭圆型有内部不动点，因此不能作为曲面无分歧覆盖的非平凡甲板变换。这把自同构的动力分类与自由覆盖作用联系起来。

% !TeX root = ../../Book3.tex
\section{圆盘与上半平面模型}
\label{b3:sec:1303}

半平面中的度量元更简单，距离公式也能直接显示到实边界和无穷远的行为。Cayley 变换保证两模型的所有几何结论一致。

\subsection{圆盘模型}
\label{b3:subsec:130301}

圆盘模型适于以内点为观察中心，测地线为正交圆弧，伪双曲球仍是欧氏圆盘。0到半径 \(r\) 点的距离是 \(\log((1+r)/(1-r))\)，所以 \(r\uparrow1\) 时对数发散。

\ProseParagraphBreak
欧氏圆心不是内在球心的欧氏位置，除非内在球心为0。给定内在半径 \(s\)，伪双曲半径为 \(\tanh(s/2)\)，再用第6章的球公式才能画出正确的欧氏圆。

\subsection{半平面模型}
\label{b3:subsec:130302}

半平面长度元为 \(|\mathrm{d}z|/\operatorname{Im}z\)。距离满足
\[
 \cosh\mathrm d_{\mathbb H}(z,w)
 =1+\frac{|z-w|^2}{2\operatorname{Im}z\operatorname{Im}w}.
\]
由半平面伪双曲分式 \(\delta=|z-w|/|z-\overline w|\)，代入 \(\cosh(2\operatorname{arctanh}\delta)=(1+\delta^2)/(1-\delta^2)\)，并用两个分母平方之差 \(4\operatorname{Im}z\operatorname{Im}w\) 即得公式。

\ProseParagraphBreak
同一竖线上 \(\mathrm d(x+\mathrm{i}y_1,x+\mathrm{i}y_2)=|\log(y_2/y_1)|\)。竖线是测地线，其他测地线是以实轴点为圆心的半圆。水平直线通常不是测地线，即使沿它可以方便地计算一条路径长度。

\ProseParagraphBreak
上半平面中的竖直线段给极简计算。连接 \(\mathrm{i}y_1\) 与 \(\mathrm{i}y_2\) 的长度是 \(|\log(y_2/y_1)|\)，因为 \(\int\mathrm{d}y/y\) 正是长度积分。水平线 \(y=c\) 的一段 Euclid 长度为 \(b-a\)，双曲长度则为 \((b-a)/c\)，但一般不是两端之间的最短路。保持同一高度的长水平段与上方弯起的正交半圆相比，会累积更长的长度；路径长度公式与距离公式分别回答这两个问题。

\ProseParagraphBreak
在半平面中，平移 \(w\mapsto w+t\) 与正伸缩 \(w\mapsto a w\) 都保持 \(|\mathrm{d}w|/\operatorname{Im}w\)。第十一章的实 Möbius 变换由这些及反演组合，并满足同样不变性。半平面到圆盘的 Cayley 变换把无穷远送到一个圆周点，因此圆盘中看起来特殊的某个边界点，可能只是换坐标后的无穷远；双曲几何没有给它额外的内部地位。

\subsection{Cayley 变换}
\label{b3:subsec:130303}

对 \(C(z)=(z-\mathrm{i})/(z+\mathrm{i})\)，有 \(|C'|=2/|z+\mathrm{i}|^2\)，以及 \(1-|C|^2=4\operatorname{Im}z/|z+\mathrm{i}|^2\)。所以
\[
 \frac{2|C'(z)|}{1-|C(z)|^2}=\frac1{\operatorname{Im}z}.
\]
这同时核对密度归一化、距离式和曲率。若圆盘密度误写为 \(1/(1-|z|^2)\)，半平面密度会变成 \(1/(2y)\)，曲率也变为 \(-4\)，后面的 Ahlfors--Schwarz 常数将整体错误。

\subsection{无穷远边界}
\label{b3:subsec:130304}

圆周或扩充实轴是模型的理想边界，不属于度量空间。两个不同边界点确定一条完整测地线；一个边界点还确定一族与边界相切的欧氏圆，在半平面中为水平线，它们称为等距边界曲线或 horocycle。它们与测地线正交，但自身不最短。

\ProseParagraphBreak
例如 \(\mathrm{i}y\) 当 \(y\downarrow0\) 或 \(y\uparrow\infty\) 都逃向理想边界，与 \(\mathrm{i}\) 的距离均趋于无穷。欧氏意义的“离原点很近”和“很远”在内在几何中都可表示同一种无限远离。

\ProseParagraphBreak
圆盘边界是理想边界，不是完备化所增添的 Cauchy 极限。向圆周径向走，距离为 \(\log((1+r)/(1-r))\to\infty\)，不可能是一条趋向有限距离点的 Cauchy 路径。比如 \(r_n=1-e^{-n}\)，相邻点的双曲距离趋近1；虽然 Euclid 差趋于0，它在双曲度量下不是 Cauchy 列。两种拓扑在内部相同，趋向边界的度量行为却不同。

\ProseParagraphBreak
以半平面无穷远为理想端点的等距平移轨道 \(y=c\) 称为水平等高曲线的模型；其曲线长度随水平参数线性增加。经 Möbius 变换，它们变为与圆盘边界相切的圆。这样的圆与测地线不同：测地圆弧正交于边界，等高曲线相切于边界。两者在图形上都可能是圆弧，判别应看边界接触角和它们承担的几何任务。

\begin{proposition}[平移与伸缩的位移长度]\label{b3:prop:13-deck-displacement}
在上半平面模型中，双曲变换 \(T(w)=aw\)、\(a>1\) 的最小位移为 \(\log a\)，由虚轴上的点达到。抛物变换 \(P(w)=w+1\) 的位移下确界为0，但没有点达到。
\end{proposition}
\begin{proof}
令 \(w=x+\mathrm{i}y\)。距离公式给
\[
 \cosh d(w,aw)=1+\frac{(a-1)^2(x^2+y^2)}{2ay^2}
 \ge\frac{a+a^{-1}}2=\cosh\log a.
\]
等号恰在 \(x=0\)。平移情形则 \(\cosh d(w,w+1)=1+1/(2y^2)\)，对任何有限 \(y\) 大于1，随 \(y\to\infty\) 趋于1。
\end{proof}
循环双曲商中，虚轴下降为一条核心闭测地线，长度 \(\log a\)；第十二章的圆环模量为 \(\pi/\log a\)，与正文核心公式一致。循环抛物商则是去心盘，绕行圈可以任意短但没有最短非平凡圈。这一差别来自甲板变换的位移是否达到正最小值，而不是来自两曲面基本群不同。

% !TeX root = ../../Book3.tex
\section{一般双曲 Riemann 曲面}
\label{b3:sec:1304}

一般双曲曲面不必有一个单值圆盘坐标，但每个小片都可以由圆盘覆盖提供坐标。甲板不变性将长度元下降，路径提升则把距离、压缩与完备性传下来。

\subsection{从万有覆盖拉回}
\label{b3:subsec:130401}

\begin{theorem}[深层覆盖模型输入]\label{b3:thm:13-cover-model-input}
每张 Riemann 曲面的全纯万有覆盖双全纯等价于球面、平面、圆盘之一。特别地，补集至少含两个点的平面区域具有圆盘万有覆盖。
\end{theorem}
\begin{proofsketch}
主线在这里独立采用统一化输入，不要求先读选读第12章。正式来源为 Ahlfors《Conformal Invariants》\cite{Ahlfors2010Conformal}定理10-3及 \S10-4--10-6（142--151页）。平面域特例由自由甲板群分类推出：球面覆盖使底面紧，不可能；平面甲板群只能为离散平移群，底面是平面、去心平面或紧环面；平面单叶映射为仿射，去心平面单叶映射在两端由 Casorati--Weierstrass 与局部逆排除本性奇点，延为一次有理映射，所以其平面像只少一个点。补集至少两点排除所有这些例外。一般统一化的势论证明作为深层来源输入，后面的度量结论逐项完整证明。
\end{proofsketch}
称圆盘覆盖的曲面为双曲曲面。给 \(p:\mathbb D\rightarrow X\)，在底坐标中定义 \(\lambda_X(p(w))=2/((1-|w|^2)|p'(w)|)\)。若两提升为 \(w,Tw\)，由 \(p\circ T=p\) 与圆盘密度不变性得相同数值，故良定义。换底坐标时按导数模变换，所以给内在光滑长度元，局部曲率为 \(-1\)。

\subsection{双曲区域}
\label{b3:subsec:130402}

对去心盘用指数半平面覆盖直接得到
\[
 \lambda_{\mathbb D^*}(z)=\frac1{|z|\log(1/|z|)}.
\]
沿径向趋于0的长度为 \(\int\mathrm{d}r/(r\log(1/r))\)，发散；穿孔虽为一个欧氏孤立点，仍在无限距离。靠近外圆周密度也发散。

\ProseParagraphBreak
对 \(A(r,R)\)，令 \(L=\log(R/r)\)，将对数带以 \(e^{\pi\mathrm{i}(w-\log r)/L}\) 映到半平面，再下降得
\[
 \lambda_A(z)=\frac{\pi}{L|z|\sin(\pi\log(|z|/r)/L)}.
\]
中间圆 \(|z|=\sqrt{rR}\) 的长度为 \(2\pi^2/L=\pi/\operatorname{mod}A\)。任意绕孔一次的闭路长度不小于此数，因为正弦不超过1且总辐角变化至少 \(2\pi\)；中间圆达到下界。

\begin{proposition}[圆环的最短核心曲线]\label{b3:prop:13-annulus-core}
设圆环的模量为 \(m>0\)。在绕孔一次的闭曲线中，双曲长度最小值为 \(\pi/m\)，由几何平均半径 \(\sqrt{rR}\) 的圆达到。绕孔 \(n\) 次的长度至少为 \(|n|\pi/m\)。
\end{proposition}
\begin{proof}
写 \(L=\log(R/r)=2\pi m\)、\(t=\log(|z|/r)\)，密度公式给
\[
 \lambda_A(z)|\mathrm{d}z|=\frac\pi{L\sin(\pi t/L)}\frac{|\mathrm{d}z|}{|z|}.
\]
沿闭路连续选辐角 \(\theta\)，\(|\mathrm{d}z|/|z|\ge|\mathrm{d}\theta|\)，且 \(\sin(\pi t/L)\le1\)。因此长度至少为 \((\pi/L)|\theta(1)-\theta(0)|=2\pi^2|n|/L\)。在 \(t=L/2\) 的圆上，正向绕行没有角方向往返，两项不等式均等号，所以达到一次绕行的最小值。多重绕行也给相应等号。
\end{proof}
模量大意味着核心闭曲线短，而连接两端的极值长度大。两项结论衡量不同曲线族，彼此并不矛盾。若内半径趋于0，模量趋于无穷，核心长度趋于0；极限去心盘没有实际达到长度0的闭曲线，短圈越来越靠近理想端部。

\subsection{共形不变性}
\label{b3:subsec:130403}

\begin{theorem}\label{b3:thm:13-hyperbolic-schwarz-pick}
双曲曲面间全纯映射 \(f:X\rightarrow Y\) 满足 \(\lambda_Y(f(z))|f'(z)|\le\lambda_X(z)\)，并压缩路径距离。
\end{theorem}
\begin{proof}
复合 \(f\circ p_X\) 从单连通圆盘出发，可用路径提升与同伦提升取全纯提升 \(F:\mathbb D\rightarrow\mathbb D\)，使 \(p_Y\circ F=f\circ p_X\)。提升在小片由全纯层逆给出，故全纯。对 \(F\) 用圆盘 Schwarz--Pick，再以覆盖导数链式法则和密度定义消去两覆盖因子，得到所述式。沿曲线积分并取下确界给距离压缩。
\end{proof}
双全纯映射及其逆各给一个不等式，因此必须等号；一般覆盖映射也是局部等距，但不同底点的最短距离可以因不同提升选择而缩短。局部等距并不要求全局距离逐对相等。

\ProseParagraphBreak
全纯映射使长度压缩时，并没有保证不同点的像不同。\(z\mapsto z^m\) 在圆盘中是压缩映射，0处导数为0；两个不同点也可能映到同一点。若导数在某一点达到 Schwarz--Pick 等号，提升后的圆盘映射才被迫是自同构；下降到曲面以后得到的是局部等距覆盖，仍可能多层。对圆盘自映射它就是自同构，对圆环间 \(z^m\) 则恰是覆盖，显示目标拓扑影响等号结论的最终表述。

\ProseParagraphBreak
把覆盖本身看作全纯映射，密度等号处处成立，这是密度下降定义的直接后果。覆盖可以将源中的不同点识别，底面两点的最短距离是所有甲板移动后的提升距离下确界，所以不要求任意固定两条提升都保持端点距离。局部长度保持与全局距离保持有区别，后者只对选好最短相对位置的提升成立。

\subsection{单调性}
\label{b3:subsec:130404}

\begin{corollary}\label{b3:cor:13-density-monotonicity}
若双曲平面区域 \(U\subset V\)，则 \(\lambda_V\le\lambda_U\)。若 \(U\ne V\)，在每个内点都严格。
\end{corollary}
\begin{proof}
对包含映射使用前一定理。若某点取等，其圆盘提升由 Schwarz--Pick 刚性为自同构。共轭后源甲板群包含于目标甲板群，底映射是相应群商之间的覆盖，特别满射；但包含映射满射只可能 \(U=V\)。故真包含没有等号点。
\end{proof}
去心盘包含于圆盘，直接公式也给 \(\lambda_{\mathbb D^*}>\lambda_{\mathbb D}\)。删点使全纯映射的活动空间更小，所以局部内在密度增大，而不是减小。

\begin{proposition}[凸域中的边界距离估计]\label{b3:prop:13-convex-boundary-distance}
设 \(U\subsetneq\mathbb C\) 为凸区域，\(d(z)=\operatorname{dist}(z,\partial U)\)。则
\[
 \frac1{d(z)}\le\lambda_U(z)\le\frac2{d(z)}.
\]
\end{proposition}
\begin{proof}
以 \(z\) 为中心、半径 \(d(z)\) 的盘包含于 \(U\)，密度单调性给上界。取最近边界点 \(q\)。凸性使经过 \(q\)、垂直于 \(z-q\) 的直线为支持线，\(U\) 落在含 \(z\) 的开半平面 \(H\) 中。可从闭凸集在 \(q\) 的支持直线存在性及内切圆的切线唯一性看出这一点；这里使用实平面的基本凸性分离事实。半平面密度在 \(z\) 为 \(1/d(z)\)，再用单调性得下界。
\end{proof}
半平面达到下界，圆盘中心达到上界，故两个常数在凸域类中不能整体改进。对于任意非凸域，上界仍由内盘包含成立，下界却不能不经新证明沿用。去心盘靠近穿孔时 \(\lambda\sim1/(|z|\log(1/|z|))\)，Euclid 边界距离为 \(|z|\)，已小于 \(1/d(z)\)，显示凸性确实参与下界。

\ProseParagraphBreak
下界也解释凸边界的无限距离。边界距离函数是1-Lipschitz，沿任意路径 \(|\mathrm{d}d|\le|\mathrm{d}z|\)，所以双曲长度至少为 \(\int|\mathrm{d}d|/d\)。若路径从固定正距离走到边界距离0，该积分至少为相应对数差并趋于无穷。这里无需假定路径径向，也无需边界可微；凸性及支持半平面已经提供整个路径上的下界。

\subsection{完备性}
\label{b3:subsec:130405}

\begin{theorem}\label{b3:thm:13-hyperbolic-complete}
双曲曲面的 Poincaré 路径度量完备。任意两点有最短测地线，且对选定提升 \(u,v\) 有
\[
 \mathrm d_X(pu,pv)=\inf_{T\in\operatorname{Deck}(p)}\mathrm d_{\mathbb D}(u,Tv).
\]
\end{theorem}
\begin{proof}
任意底路径从 \(u\) 提升，终点是某个 \(Tv\)，长度保持，所以底距离至少右端；任意圆盘最短路投影给底路径，反向不等式成立。纤维是闭离散集，圆盘闭有界球紧，故在给定有限距离内只有有限个 \(Tv\)；下确界由其中一点达到，最短路投影得到底测地线。

设底空间有 Cauchy 列，抽取子列 \(x_n\) 满足相邻距离小于 \(2^{-n}\)，选连接路径长度小于 \(2^{1-n}\)。从选定提升开始逐次提升这些路径，得到圆盘点 \(u_n\)。尾部连接长度和趋于0，使 \(u_n\) 在圆盘为 Cauchy，完备性给极限 \(u\)。投影后 \(x_n\rightarrow p(u)\)，原 Cauchy 列也趋于同一点。因此底度量完备。
\end{proof}
在环域中不同绕行类可能给不同测地线，所以唯一性不再像圆盘那样自动成立。去心盘的绕孔闭路长度可趋于0，却没有长度0的闭测地线；完备性不能误解为每个非平凡闭路类都有长度极小代表。

\ProseParagraphBreak
覆盖距离公式还能给局部没有自重叠的半径。对 \(x=p(w)\)，令
\[
 r_x=\tfrac12\inf_{T\ne1}d_{\mathbb D}(w,Tw),
\]
若甲板群平凡便取无穷。自由离散作用和闭球紧使这个下确界为正：否则非平凡轨道点将趋于 \(w\)，违反覆盖均匀邻域的分离。小于 \(r_x\) 的圆盘球投影单射；若两点被 \(T\) 识别，三角不等式会使 \(d(w,Tw)<2r_x\)，矛盾。这给曲面上围绕 \(x\) 的一张无重叠等距坐标图。

\ProseParagraphBreak
对于有限圆环，核心附近的这种半径由核心长度控制；对去心盘，向穿孔端行进时绕行圈的长度趋于0，所以局部单射半径也可能趋于0。曲面仍完整，因为通向穿孔的径向长度 \(\int\mathrm{d}r/[r\log(1/r)]\) 发散。越来越短的绕行闭路与越来越远的端部可以同时存在，完整度量并没有要求全曲面的局部单射半径有统一正下界。

\begin{proposition}[商曲面的闭球与最短路]\label{b3:prop:13-quotient-geodesics}
圆盘型曲面的有限半径闭球紧，每两点间都有长度等于距离的路径；最短路径不必唯一。
\end{proposition}
\begin{proof}
固定 \(x=p(w)\)。覆盖距离公式及离散纤维说明，对任意 \(y\)，可在其纤维中找到 \(v\) 使 \(d_X(x,y)=d_{\mathbb D}(w,v)\)：若距离下确界有限，取一个有界距离序列，圆盘闭球紧而纤维闭离散，故其中只有有限个点，下确界达到。因此半径 \(R\) 的闭球恰是 \(p(\overline B_{\mathbb D}(w,R))\)，是紧集。把 \(w\) 到最佳提升 \(v\) 的圆盘测地段投影，其长度保持并等于底面距离，故是最短路。
\end{proof}
有限圆环的核心圆上两个相对点有两条最短路径，分别沿核心走半圈。任意连接路径的连续辐角变化为 \(\pi+2\pi n\)，绝对值至少 \(\pi\)；核心长度证明中的角积分下界说明半圈达到最小值。圆盘自身的最短路唯一，商以后却因不同甲板相对位置取得同样最小值而不唯一。局部测地线仍由圆盘模型决定，全局唯一性不随覆盖自动下降。


\begin{chaptersummary}
Poincaré 度量从局部密度出发，经路径积分与长度下确界得到两点距离。在圆盘上，自同构把一端移到0，径向积分给出 \(2\operatorname{arctanh}\delta\)；等号条件随后给出唯一最短段。边界的无限距离使圆盘完备，而有限半径闭球仍是内部紧集。半平面模型把同一几何写成 \(|\mathrm dw|/\operatorname{Im}w\)，所有距离、面积与曲率公式都要保持同一个归一化。

\ProseParagraphBreak
圆盘覆盖下降的是长度元，路径提升下降的是距离比较。底曲面距离等于不同甲板相对位置下提升距离的最小值，因此最短路存在，却可能有多条。全纯映射压缩长度和距离；微分等号迫使映射为覆盖，只有在进一步知道双射时才能断言它保持全部两点距离。区域包含给出密度单调性，凸性则让边界距离提供统一的双边估计。

\ProseParagraphBreak
去心盘和圆环使局部与全局的差别具体可算。趋向穿孔的路径需无限长度，附近的绕行圈却可以很短，尖端尾部还可以只有有限面积。单射半径描述局部不重叠的范围，不由完备性提供统一正下界。紧闭球与全纯压缩结合后，固定基点像的映射族便不能向端部逃走；这给出曲面值极限，而不是把所有边界退化都纳入同一种收敛。
\end{chaptersummary}

\BookExercises

% 八题均为自拟或对原自拟题的实质改编；解答仅用已建立的距离、覆盖和换元工具，不附会未实读的外部题源。
\Needspace{6\baselineskip}
\begin{exercise}[夹角与观察尺度]\label{b3:ex:13-1}
在曲率为 \(-1\) 的圆盘模型中，取从0出发、夹角为 \(\theta\in(0,\pi]\) 的两条径向测地线。两端点到0的双曲距离都为 \(R>0\)，端点之间的距离记为 \(c(R,\theta)\)。证明
\[
 \sinh\frac{c(R,\theta)}2=\sinh R\cdot\sin\frac\theta2.
\]
对固定 \(\theta\)，求 \(R\downarrow0\) 和 \(R\rightarrow\infty\) 时的首项，并写出直角情形的结果。再固定 \(A>0\)，令夹角随距离缩小为 \(\theta_R=Ae^{-R}\)，求 \(c(R,\theta_R)\) 的极限；这里只考虑使 \(\theta_R\le\pi\) 的充分大 \(R\)。
\end{exercise}
\begin{solution}
由\cref{b3:prop:13-cosh-distance}，
\[
 \cosh c=\cosh^2R-\sinh^2R\cos\theta
 =1+2\sinh^2R\cdot\sin^2\frac\theta2.
\]
再用 \(\cosh c-1=2\sinh^2(c/2)\)，且两边的平方根均为正，便得题设恒等式。因此
\[
 c(R,\theta)=2\operatorname{arsinh}\left(\sinh R\cdot\sin\frac\theta2\right).
\]
固定 \(\theta>0\)，由 \(\sinh R\sim R\)、\(\operatorname{arsinh}u\sim u\)，小半径时有
\[
 c(R,\theta)\sim2R\sin\frac\theta2\quad(R\downarrow0).
\]
这正是欧氏等腰三角形底边的首项。

当 \(R\rightarrow\infty\) 时，\(\sinh R=e^R(1+o(1))/2\)，而 \(\operatorname{arsinh}u=\log(2u)+o(1)\) 对 \(u\rightarrow\infty\) 成立，所以
\[
 c(R,\theta)=2R+2\log\sin\frac\theta2+o(1).
\]
直角情形给 \(\cosh c=\cosh^2R\)，以及 \(c\sim\sqrt2R\)、\(c=2R-\log2+o(1)\)。

若 \(\theta_R=Ae^{-R}\)，则角度与半径的变化必须同时保留：
\[
 \sinh R\cdot\sin\frac{Ae^{-R}}2\longrightarrow\frac A4.
\]
故 \(c(R,\theta_R)\rightarrow2\operatorname{arsinh}(A/4)\)。两端点各自逃向理想边界，彼此距离却趋于有限正数。固定夹角时所得的大半径展开并不对所有趋零的夹角一致有效。
\end{solution}

\Needspace{6\baselineskip}
\begin{exercise}[先上升再横行的路径]\label{b3:ex:13-2}
在上半平面 \(\mathbb H\) 中取长度元 \(|\mathrm dw|/\operatorname{Im}w\)。连接 \(\mathrm i\) 与 \(a+\mathrm i\)、\(a>0\)，考虑三段路径：先沿竖线升到 \(\mathrm i h\)，再水平走到 \(a+\mathrm i h\)，最后竖直下降，其中 \(h\ge1\)。求该类路径的最短长度 \(L_*(a)\) 及最优高度，并与真正的双曲距离 \(D(a)\) 比较。证明对所有 \(a>0\) 都有
\[
 0<L_*(a)-D(a)<2-2\log2,
\]
且右侧常数不能减小。解释这如何改进始终沿高度1水平前进的路径。
\end{exercise}
\begin{solution}
两段竖直路径各长 \(\log h\)，水平段长 \(a/h\)，因此
\[
 L_h(a)=2\log h+\frac ah,\quad
 \frac{\partial L_h}{\partial h}=\frac{2h-a}{h^2}.
\]
若 \(0<a\le2\)，导数在 \(h\ge1\) 上非负，最优高度为1；若 \(a>2\)，先减后增，最优高度为 \(a/2\)。所以
\[
 L_*(a)=
 \begin{cases}
 a,&0<a\le2,\\
 2\log(a/2)+2,&a\ge2.
 \end{cases}
\]
半平面距离公式给
\[
 \cosh D(a)=1+\frac{a^2}{2},\quad
 D(a)=2\operatorname{arsinh}(a/2).
\]
记差值为 \(E(a)=L_*(a)-D(a)\)。当 \(0<a<2\) 时，
\[
 E'(a)=1-\frac1{\sqrt{1+a^2/4}}>0.
\]
当 \(a>2\) 时，
\[
 E'(a)=\frac2a-\frac2{\sqrt{a^2+4}}>0.
\]
两段在2处连续，且 \(E(a)\rightarrow0\) 当 \(a\downarrow0\)。故 \(E(a)>0\) 并严格递增。又由 \(D(a)=2\log a+o(1)\)，
\[
 E(a)\longrightarrow2-2\log2\quad(a\rightarrow\infty).
\]
因此题设上下界成立，上界为差值的上确界，不能减小。

高度1的水平路径长为 \(a\)，与真正距离的差趋于无穷。允许先上升后，较大的高度使横向密度降低，竖直路程却只付出对数长度。三段路径虽然仍不是真正的测地圆弧，却在任意尺度上只多花一个统一有界的长度。
\end{solution}

\Needspace{6\baselineskip}
\begin{exercise}[相邻距离与 Cauchy 条件]\label{b3:ex:13-3}
在单位圆盘的实直径上，令
\[
 r_n\coloneqq1-e^{-n}\quad(n\ge1),\quad
 s_n\coloneqq1-\frac1n\quad(n\ge2).
\]
证明两列都是欧氏 Cauchy 列，且都没有双曲 Cauchy 子列。进一步证明：\(r_n\) 的任意两个不同项之间的双曲距离大于1，而 \(s_n\) 的相邻项双曲距离趋于0；直接计算 \(\mathrm d_{\mathbb D}(s_n,s_{2n})\) 的极限。最后用 \(r_n\) 的相邻项说明，从欧氏圆盘到 Poincaré 圆盘的恒等映射虽连续，却不一致连续。
\end{exercise}
\begin{solution}
两列在欧氏意义下都趋于1，所以都是欧氏 Cauchy 列。实直径上的径向弧长参数为
\[
 S(t)\coloneqq\log\frac{1+t}{1-t}\quad(0\le t<1),
\]
距离公式给 \(\mathrm d_{\mathbb D}(u,v)=|S(u)-S(v)|\)。代入 \(r_n\)，得到
\[
 S(r_n)=n+\log(2-e^{-n}).
\]
其中第二项严格递增，故 \(m>n\) 时
\[
 \mathrm d_{\mathbb D}(r_m,r_n)>m-n\ge1.
\]
任何子列都不能满足 Cauchy 条件。

对于 \(s_n\)，有 \(S(s_n)=\log(2n-1)\)，因此
\[
 \mathrm d_{\mathbb D}(s_n,s_{n+1})
 =\log\frac{2n+1}{2n-1}\longrightarrow0,
\]
但
\[
 \mathrm d_{\mathbb D}(s_n,s_{2n})
 =\log\frac{4n-1}{2n-1}\longrightarrow\log2>0.
\]
所以相邻项越来越近不能推出 Cauchy 条件。事实上，\(\mathrm d_{\mathbb D}(0,s_n)=\log(2n-1)\rightarrow\infty\)，任一子列到0的距离也无界，而每个 Cauchy 列必有界，故这里同样没有 Cauchy 子列。

最后，\(|r_{n+1}-r_n|=(1-e^{-1})e^{-n}\rightarrow0\)，对应双曲距离却始终大于1。这与一致连续性要求矛盾。普通连续性仍成立，因为在每个固定内点附近，Poincaré 密度有有限的正上下界；所缺的是同时控制整个圆盘的共同界。
\end{solution}

\Needspace{6\baselineskip}
\begin{exercise}[共同像界与复合顺序]\label{b3:ex:13-4}
固定 \(0<r<1\)，设全纯映射 \(f_j:\mathbb D\rightarrow r\mathbb D\)、\(j\ge1\)。证明每个 \(f_j\) 关于 Poincaré 距离都是 Lipschitz 常数至多为 \(r\) 的映射。再证明
\[
 \Phi_n\coloneqq f_1\circ f_2\circ\cdots\circ f_n
\]
在整个 \(\mathbb D\) 上关于双曲距离一致趋于某个盘内常值。给出只依赖 \(r,n\) 的误差界。若把复合顺序改为 \(\Psi_n=f_n\circ\cdots\circ f_1\)，结论是否仍成立？请用全部非恒常的仿射映射构造反例。
\end{exercise}
\begin{solution}
对任一 \(f_j\) 写 \(f_j=rg_j\)，其中 \(g_j:\mathbb D\rightarrow\mathbb D\) 全纯。Schwarz--Pick 导数式给
\[
 \frac{|f_j'(z)|(1-|z|^2)}{1-|f_j(z)|^2}
 \le r\frac{1-|g_j(z)|^2}{1-r^2|g_j(z)|^2}\le r.
\]
左边就是目标密度乘以导数模后，与源密度之比。沿任意路径积分并取长度下确界，得
\[
 \mathrm d_{\mathbb D}(f_j(z),f_j(w))\le r\,\mathrm d_{\mathbb D}(z,w).
\]

令 \(R\coloneqq\log((1+r)/(1-r))\)。任一 \(r\mathbb D\) 中的点到0距离小于 \(R\)，所以该集合的双曲直径至多 \(2R\)。最内层 \(f_n\) 先把整个圆盘送入这个有界集合，其余 \(n-1\) 次复合各至多压缩 \(r\) 倍，故
\[
 \operatorname{diam}_{\mathrm d_{\mathbb D}}\Phi_n(\mathbb D)
 \le2Rr^{n-1}.
\]
当 \(m\ge n\) 时，\(\Phi_m(0)\) 也属于 \(\Phi_n(\mathbb D)\)，因为把 \(f_{n+1}\circ\cdots\circ f_m(0)\) 看作 \(\Phi_n\) 的输入即可。因此 \(\Phi_n(0)\) 为双曲 Cauchy 列，由圆盘完备性有盘内极限 \(a\)。固定 \(z\in\mathbb D\)，让 \(m\rightarrow\infty\)，同一估计给
\[
 \mathrm d_{\mathbb D}(\Phi_n(z),a)\le2Rr^{n-1}.
\]
右边与 \(z\) 无关，故为所求全盘一致误差界。

反向顺序不保证收敛。令 \(q=r/3\)，并取
\[
 f_{2k-1}(z)=qz+q,\quad f_{2k}(z)=qz-q\quad(k\ge1).
\]
这些函数均非恒常，且 \(|f_j(z)|<2q<r\)。令 \(x_n=\Psi_n(0)\)、\(x_0=0\)，便有
\[
 x_{2k}=q^2x_{2k-2}+q^2-q.
\]
由于 \(0<q<1\)，这个递推给
\[
 x_{2k}\longrightarrow-\frac q{1+q},\quad
 x_{2k-1}=qx_{2k-2}+q\longrightarrow\frac q{1+q}.
\]
两极限不同，所以连0处也不收敛。共同压缩控制不同初值之间的距离，却不能阻止每次改变最外层映射时，像的中心来回移动。
\end{solution}

\Needspace{6\baselineskip}
\begin{exercise}[理想端点的极限距离]\label{b3:ex:13-5}
在上半平面的 Poincaré 距离下，对 \(w=x+\mathrm i y\in\mathbb H\) 研究
\[
 b_\infty(w)\coloneqq
 \lim_{t\rightarrow\infty}\bigl(\mathrm d_{\mathbb H}(w,\mathrm i e^t)-t\bigr).
\]
证明极限存在，求出其值，并证明 \(b_\infty\) 关于双曲距离是1-Lipschitz函数。对任意 \(a\in\mathbb R\)，同样求
\[
 b_a(w)\coloneqq
 \lim_{t\rightarrow\infty}\bigl(\mathrm d_{\mathbb H}(w,a+\mathrm i e^{-t})-t\bigr).
\]
描述两类函数的全部等值线，说明它们与正文中的极限圆有何关系。
\end{exercise}
\begin{solution}
半平面距离公式给
\[
 \cosh\mathrm d_{\mathbb H}(w,\mathrm i e^t)
 =\frac{x^2+y^2+e^{2t}}{2ye^t}.
\]
右边趋于无穷，使用 \(\operatorname{arcosh}u=\log(2u)+o(1)\)，得到
\[
 \mathrm d_{\mathbb H}(w,\mathrm i e^t)
 =t-\log y+o(1).
\]
故 \(b_\infty(w)=-\log y\)。这里减去的 \(t\) 正好是 \(\mathrm i\) 到 \(\mathrm i e^t\) 的距离；两个发散距离之差留下有限的位置信息。

对任意 \(w,v\in\mathbb H\)，三角不等式给
\[
 \left|\mathrm d_{\mathbb H}(w,\mathrm i e^t)
       -\mathrm d_{\mathbb H}(v,\mathrm i e^t)\right|
 \le\mathrm d_{\mathbb H}(w,v).
\]
令 \(t\rightarrow\infty\)，即得
\( |b_\infty(w)-b_\infty(v)|\le\mathrm d_{\mathbb H}(w,v) \)。

半平面自同构 \(M_a(w)=-1/(w-a)\) 把 \(a+\mathrm i e^{-t}\) 送到 \(\mathrm i e^t\)，且
\[
 \operatorname{Im}M_a(w)=\frac{y}{|w-a|^2}.
\]
由等距性，\(b_a(w)=b_\infty(M_a(w))\)，所以
\[
 b_a(w)=\log\frac{|w-a|^2}{y}.
\]
复合等距映射也说明 \(b_a\) 是1-Lipschitz函数。

\(b_\infty=c\) 的等值线为水平直线 \(y=e^{-c}\)。而 \(b_a=c\) 等价于
\[
 (x-a)^2+y^2=e^c y,
\]
也就是以 \(a+\mathrm i e^c/2\) 为圆心、\(e^c/2\) 为半径的欧氏圆删去切点 \(a\)。它们正是指向相应理想端点的极限圆。其定义使用趋向无穷远时的距离差，不把理想端点当作曲面中的有限距离点。
\end{solution}

\Needspace{6\baselineskip}
\begin{exercise}[到测地线的最近点]\label{b3:ex:13-6}
在上半平面模型中，令 \(L=\{\,\mathrm i t\mid t>0\,\}\) 为竖直测地线。对 \(w=x+\mathrm i y\in\mathbb H\)，证明 \(L\) 上有唯一最近点，并求最近点 \(P(w)\) 与距离 \(\mathrm d_{\mathbb H}(w,L)\)。证明投影 \(P:\mathbb H\rightarrow L\subset\mathbb H\) 不增加任意两点的双曲距离，却不是全纯映射。对于 \(w\notin L\)，说明连接 \(w\) 与 \(P(w)\) 的测地段为何垂直于 \(L\)。
\end{exercise}
\begin{solution}
由距离公式，对 \(t>0\) 有
\[
 \cosh\mathrm d_{\mathbb H}(w,\mathrm i t)
 =\frac{x^2+y^2+t^2}{2yt}
 =\frac1{2y}\left(t+\frac{|w|^2}{t}\right).
\]
\(\cosh\) 在非负实轴严格递增，而括号内的正数之和由算术几何平均不等式恰在 \(t=|w|\) 时取唯一最小值。因此
\[
 P(w)=\mathrm i|w|,\quad
 \mathrm d_{\mathbb H}(w,L)
 =\operatorname{arcosh}\frac{|w|}{y}
 =\operatorname{arsinh}\frac{|x|}{y}.
\]
最后一个等式由 \(|w|^2=x^2+y^2\) 及两个反函数的非负取值确定。

对上半平面内任意分段 \(C^1\) 路径 \(\gamma\)，在可微点有
\[
 \left|\frac{\mathrm d}{\mathrm dt}\log|\gamma(t)|\right|
 \le\frac{|\gamma'(t)|}{|\gamma(t)|}
 \le\frac{|\gamma'(t)|}{\operatorname{Im}\gamma(t)}.
\]
沿连接 \(w,v\) 的路径积分并取下确界，得到
\[
 |\log|w|-\log|v||\le\mathrm d_{\mathbb H}(w,v).
\]
左边恰为 \(P(w),P(v)\) 在同一竖直测地线上的距离，所以 \(P\) 不增距。它非恒常，却把整个上半平面映入没有内部的虚轴；若它全纯，就与开映射定理矛盾。因此不增距条件本身不能反推全纯性。

当 \(w\notin L\) 时，以0为圆心、\(|w|\) 为半径的上半圆同时经过 \(w\) 与 \(\mathrm i|w|\)，并与实轴正交，所以包含两点间的测地段。它在 \(\mathrm i|w|\) 的切线水平，与 \(L\) 垂直。共形度量保持夹角，故这个欧氏直角也是双曲几何中的直角。
\end{solution}

\Needspace{6\baselineskip}
\begin{exercise}[单射半径的变化与有限面积]\label{b3:ex:13-7}
设 \(X\) 是连通双曲 Riemann 曲面，\(p:\mathbb D\rightarrow X\) 为万有覆盖，甲板群 \(G\) 非平凡。采用正文的单射半径
\[
 r_x\coloneqq\frac12\inf_{T\in G,\ T\ne1}
 \mathrm d_{\mathbb D}(u,Tu),\quad p(u)=x.
\]
证明 \(|r_x-r_y|\le\mathrm d_X(x,y)\)，从而每个紧集上的单射半径有统一正下界。再证明：若 \(X\) 的双曲总面积有限，且存在 \(r_0>0\) 使所有 \(r_x\ge r_0\)，则 \(X\) 必紧。逆向表述这一结论，说明有限面积的非紧曲面为何必须有单射半径趋于0的点列。
\end{exercise}
\begin{solution}
固定 \(x,y\in X\)，先取 \(p(u)=x\)，再由\cref{b3:thm:13-hyperbolic-complete}中的最小值性质选取 \(p(v)=y\)，使
\[
 \mathrm d_{\mathbb D}(u,v)=\mathrm d_X(x,y).
\]
对每个非恒等 \(T\in G\)，三角不等式和 \(T\) 的等距性给
\[
 \mathrm d_{\mathbb D}(u,Tu)
 \le2\mathrm d_{\mathbb D}(u,v)+\mathrm d_{\mathbb D}(v,Tv).
\]
对 \(T\ne1\) 取下确界并除以2，得到 \(r_x\le\mathrm d_X(x,y)+r_y\)。交换 \(x,y\) 得反向估计，故 \(r_x\) 是1-Lipschitz函数。正文已证明各点 \(r_x>0\)，所以在非空紧集上，连续函数 \(r_x\) 的最小值为正。

现在假设所有 \(r_x\ge r_0>0\)，固定 \(0<s<r_0/2\)。正文的覆盖球比较说明，每个 \(B_X(x,s)\) 都由圆盘中的半径 \(s\) 球单射、等距地投影而来。覆盖保持长度元，也保持面积元，因此每个这样的球都有相同的正面积
\[
 A_s=2\pi(\cosh s-1)>0.
\]
若 \(X\) 不紧，利用\cref{b3:prop:13-quotient-geodesics}的紧闭球性质，可以递归选取 \(x_1,x_2,\ldots\)，使任意两点间距离大于 \(2s\)：已选有限点以后，它们的半径 \(2s\) 闭球之并紧，不可能覆盖非紧的 \(X\)，故仍可在并集之外选下一点。

于是各个 \(B_X(x_j,s)\) 两两不交。对任意正整数 \(N\)，总面积至少为 \(NA_s\)，与总面积有限矛盾。所以 \(X\) 必紧。

因此，若 \(X\) 非紧而总面积有限，就不可能有统一的正单射半径下界，可选 \(x_n\) 使 \(r_{x_n}<1/n\)。这样的点列还必须离开每个固定紧集，因为前半题已给出紧集上的正下界。完备性允许这些端部存在，但有限面积迫使其局部不重叠尺度越来越小。
\end{solution}

\Needspace{6\baselineskip}
\begin{exercise}[面积、覆盖张数与紧曲面自映射]\label{b3:ex:13-8}
设 \(p:X\rightarrow Y\) 为连通双曲 Riemann 曲面之间的全纯覆盖，两曲面都取曲率为 \(-1\) 的 Poincaré 面积。证明有限 \(N\) 张覆盖满足
\[
 \operatorname{Area}(X)=N\operatorname{Area}(Y),
\]
而无限张覆盖的源总面积必为无穷。由此证明有限总面积曲面的每个全纯自覆盖都是双全纯映射，并解释去心盘上的幂覆盖为何不构成反例。最后设 \(X\) 紧，证明每个非恒常全纯自映射 \(f:X\rightarrow X\) 都是双全纯映射；此时不预先假设 \(f\) 无分歧，也不使用 Riemann--Hurwitz 公式。
\end{exercise}
\begin{solution}
覆盖是局部等距映射，故在均匀覆盖坐标中保持面积。为避免把重叠坐标片重复计数，取覆盖 \(Y\) 的可数个均匀覆盖坐标邻域 \(U_j\)，再作两两不交的 Borel 分解
\[
 A_1=U_1,\quad A_j=U_j\setminus\bigcup_{k<j}U_k\quad(j\ge2).
\]
可数覆盖来自曲面的第二可数性。若覆盖有 \(N<\infty\) 张，则每个 \(A_j\) 在 \(U_j\) 的各层中各有一份原像，每份面积等于 \(\operatorname{Area}(A_j)\)。因此
\[
 \operatorname{Area}(p^{-1}(A_j))=N\operatorname{Area}(A_j).
\]
对互不相交的 \(A_j\) 求和，即得总面积公式，面积无穷时也由非负可数可加性成立。

若覆盖有无限张，只需选一个相对紧、正面积的均匀覆盖小盘 \(U\)。它有任意多个互不相交的原像层，各层面积都等于同一个正数 \(\operatorname{Area}(U)\)，故 \(X\) 总面积为无穷。若 \(p:X\rightarrow X\) 且 \(0<\operatorname{Area}(X)<\infty\)，先排除无限张，再由面积公式得到 \(N=1\)，所以 \(p\) 双全纯。

去心盘的总面积并不有限。写 \(t=\log(1/r)\)，其面积积分为
\[
 2\pi\int_0^1\frac{\mathrm dr}{r\log^2(1/r)}
 =2\pi\int_0^\infty\frac{\mathrm dt}{t^2}=\infty.
\]
发散发生在外边界 \(r\uparrow1\)，不与穿孔尖端的尾部面积有限相矛盾。因此 \(z\mapsto z^N\)、\(N\ge2\)，可以是去心盘的多张自覆盖。

现设 \(X\) 紧且 \(f\) 非恒常。开映射定理使 \(f(X)\) 开，紧性使它闭，连通性因而给 \(f(X)=X\)。在局部坐标中，临界点是导数的零点；非恒常性和恒等原理使它们孤立，紧性又使临界点集 \(C\) 有限。每个纤维同样是紧的离散集，因而有限。于是临界值集 \(B=f(C)\) 及其全部原像 \(S=f^{-1}(B)\) 都有限。

连通曲面删去有限点仍连通：给定两个未删点，先在曲面上用路径连接，再在各被删点周围的小坐标圆盘内把经过该点的部分绕开。限制映射
\[
 f_0:X\setminus S\longrightarrow X\setminus B
\]
于是有连通源和目标，仍满射、适当且没有临界点。适当性是因为目标中任一紧集的原像在紧曲面 \(X\) 中闭，且不交 \(S\)。由\cref{b3:prop:11-proper-unbranched}，\(f_0\) 为有限覆盖，其张数 \(N\ge1\) 在连通目标上恒定。

记 \(\mathrm dA_X\) 为 Poincaré 面积元，并在源、目标的局部坐标中写
\[
 J_f(z)\coloneqq
 \left(\frac{\lambda_X(f(z))\cdot|f'(z)|}{\lambda_X(z)}\right)^2.
\]
坐标变换的导数因子抵消，故 \(J_f\) 是全局连续函数；\cref{b3:thm:13-hyperbolic-schwarz-pick}给 \(0\le J_f\le1\)。在 \(X\setminus B\) 的均匀覆盖小片上使用实变量换元，再按本题开头的 Borel 分解求和，得到
\[
 \int_X J_f\,\mathrm dA_X=N\operatorname{Area}(X).
\]
这里使用的是局部 Jacobian \(|f'|^2\)，并未预先认定 \(f\) 等距；有限集 \(B,S\) 的面积为零，所以不影响等式。

紧曲面的光滑正面积元具有有限正总面积，于是
\[
 N\operatorname{Area}(X)=\int_X J_f\,\mathrm dA_X
 \le\operatorname{Area}(X)
\]
迫使 \(N=1\) 及 \(\int_X(1-J_f)\,\mathrm dA_X=0\)。非负连续函数 \(1-J_f\) 若在某点为正，便在一个正面积邻域内有正下界，使积分为正，矛盾。因此 \(J_f\equiv1\)。导数密度处处取等，\cref{b3:thm:13-hyperbolic-schwarz-pick}说明 \(f\) 是全纯覆盖；它在正则值上的张数已经为1，故为一张覆盖，即双全纯映射。
\end{solution}
